% Part: incompleteness
% Chapter: theories-computability
% Section: complete-decidable

\documentclass[../../../include/open-logic-section]{subfiles}

\begin{document}

\olfileid{inc}{tcp}{cdc}

\olsection{可公理化的完备理论是可判定的}

如果对于每个语句 $!A$，$!A$ 或 $\lnot !A$ 可证，则称理论是{\em 完备的}。

\begin{lem}
设理论 $\Th{T}$ 是完备的且可公理化的，则 $\Th{T}$ 是可判定的。
\end{lem}

\begin{proof}
假设 $\Th{T}$ 是完备的，且 $A$ 是可计算的公理集合。
如果 $\Th{T}$ 是不一致的，则它显然是可计算的。（算法：“直接回答是。”）
因此我们可以假定 $\Th{T}$ 也是一致的。

要判定语句 $!A$ 是否属于~$\Th{T}$，
同时搜索~$!A$ 与~$\lnot !A$ 从~$\Th{T}$ 的推导。
由于 $\Th{T}$ 是完备的，总能找到其中之一；
又由于 $\Th{T}$ 是一致的，如果找到了~$\lnot !A$ 的推导，
则不存在~$!A$ 的推导。

换句话说，我们已经知道 $\Th{T}$ 是 c.e.；
所以根据之前证明过的一个定理，只需证明~$\Th{T}$ 的补集也是 c.e. 即可。
而公式~$!A$ 属于~$\Th{\bar T}$ 当且仅当 $\lnot !A$ 属于~$\Th{T}$；
因此 $\Th{\bar T} \leq_m
\Th{T}$。
\end{proof}

\end{document}
